\subsubsection*{Orientation}

\begin{remark*}[Section map]
This section begins with Boolean functions and formula representation, then
introduces connective bases and definability of connectives. It uses normal
forms to prove completeness of the standard connectives and then reduces the
primitive vocabulary further. The final topics show that NAND and NOR each form
a one-connective complete basis, while some familiar connective sets fail to be
functionally complete.
\end{remark*}

\begin{remark*}[Connection with normal forms]
Normal forms explain why functional completeness is possible. Canonical DNF and
canonical CNF show that the truth behavior of a formula can be reconstructed
from finitely many truth-table rows. Once those canonical constructions are
available, proving that a connective set is functionally complete amounts to
showing that the connectives needed for those normal forms can be defined from
that set.
\end{remark*}
